\section{The problem the serving stack actually solves}
\label{sec:orientation}
A model checkpoint describes a computation. A serving system must decide when,
where, and with what retained state to perform that computation for many users
at once. Two prompts of the same length can have very different service costs:
one may reuse a cached prefix on an idle GPU; another may need a full prefill,
a remote transfer, or admission behind a long queue. A fleet with fast kernels
can still deliver poor latency if it repeatedly recomputes old contexts or
routes requests to an overloaded cache-rich worker.

An agent makes the problem more temporal. A coding agent may generate a command,
wait for a test suite, append its output, and ask the model again. During the
wait, the GPU sees no active decode for that session. Nevertheless, retaining
some state may save work when the tool finishes. Retaining everything can also
prevent other sessions from making progress. The decision requires the cost of
keeping state, the cost of losing it, and the likely time until reuse.

\subsection{A map of responsibilities}
Figure~\ref{fig:stack} distinguishes six layers. The application determines
semantics and dependencies. A gateway enforces request limits and routes work.
An inference engine schedules token computation and manages active state.
Kernels implement tensor operations. A cache service extends where reusable
state can live. A tool runtime executes non-LLM work. A single product can span
several layers; the boundaries describe responsibilities, not mandatory processes.

\fig{stack}{0.82}{A conceptual serving stack. Arrows denote request or data exchange,
not a mandated implementation. ATFM primarily concerns the information and policy
between applications and serving; it does not replace GPU kernels.}{stack}

\begin{table}[htbp]\centering\small
\caption{Questions that belong to different layers.}
\begin{tabularx}{\linewidth}{@{}lXX@{}}\toprule
Layer & Main question & Typical cost of a wrong decision\\\midrule
Kernel & How do these tensors execute efficiently? & Extra device IO or compute\\
Engine & Which token work fits this iteration? & Stalls, preemption, wasted capacity\\
Cache & Which reusable state should exist where? & Recompute, transfer, occupied memory\\
Router & Which worker should receive this request? & Queueing or lost locality\\
Admission & Should this work enter now? & Overload or unnecessary delay\\
Agent control & What is likely to arrive after a tool? & Missed preparation or wasted reservation\\\bottomrule
\end{tabularx}\end{table}

\subsection{How to read this paper}
Readers new to serving should start with the equations and worked memory example
in Section~\ref{sec:fundamentals}, then the engine and cache sections. Readers
already operating vLLM can focus on cache movement, agent scheduling, and the
ATFM assessment. The appendix provides notation, terminology, a diagnosis guide,
and a reading sequence. Equations are models with stated assumptions, not
universal performance predictors.

The survey is broad but deliberately bounded. It emphasizes autoregressive
transformer inference, open serving infrastructure, and systems relevant to
agent traffic. Training, model-quality leaderboards, mobile inference, and every
proprietary platform are outside scope. ``State of the art'' means the relevant
research frontier and actively documented mechanisms, not one universally
fastest system. Published comparisons differ in hardware, model, workload,
cache warmness, implementation generation, and latency constraints. Their
headline speedups cannot be put in one league table.

\subsection{Three kinds of evidence}
\textbf{Established mechanism} means a cited paper or implementation documents
how a technique works. \textbf{Analytical illustration} means a derivation or a
synthetic example in this paper, with explicit numbers and assumptions.
\textbf{ATFM evidence} means a result recorded in this repository. A source's
reported benchmark is evidence for that source's configuration, not an ATFM
result. Recent preprints are marked as such; living documentation has an access
date. We use primary papers and project documentation, and avoid reproducing
vendor peak-throughput rankings.

The central hypothesis is intentionally narrower than ``agents need better
serving'': conditional information about a paused session can improve a future
resource decision enough to pay for obtaining and acting on that information.
A forecast with lower statistical error is not sufficient. The decision must
change, complete before it matters, and improve a user-relevant outcome.
